\documentclass[12pt,a4paper]{article}
% Preambule commun aux comptes rendus CR_*.tex.

\usepackage{array}
\newcolumntype{C}[1]{>{\centering\arraybackslash}m{#1}}
\usepackage{multirow}
\usepackage{geometry}
\geometry{a4paper, margin=1in}
\usepackage{graphicx}
\usepackage{float}
\usepackage{tikz}
\usetikzlibrary{arrows.meta,positioning}
\usepackage{pgfgantt}
\usepackage{pdflscape}
\usepackage{enumitem}
\usepackage{booktabs}
\usepackage{longtable}
\usepackage{ragged2e}
\usepackage{tabularx}
\usepackage{xparse}
\usepackage{ltablex}
\usepackage{xltabular}
\usepackage{subcaption}

\newcolumntype{Y}{>{\RaggedRight\arraybackslash}X}

\newcommand{\StyledTableSetup}{%
  \footnotesize
  \renewcommand{\arraystretch}{1.2}%
  \setlength{\tabcolsep}{4pt}%
}


% ============================================================
% IHEDN COVER TEMPLATE (XeLaTeX)
% ============================================================
% Compilation (from project root):
%   latexmk -pdf main.tex
%   LATEXMK_SHELL_ESCAPE=0 latexmk -pdf main.tex
%
% Optional environment controls from latexmkrc:
%   LATEXMK_ENGINE=xelatex        (default/enforced)
%   LATEXMK_SHELL_ESCAPE=1|0      (default 1, needed for SVG conversion)
%
% Required package for this template:
%   \usepackage{ihedn-cover}
%
% Note: ihedn-cover already loads needed internal packages
% (fontspec, graphicx, tikz, ragged2e, optional svg).
% ============================================================

% ----------------------------------------------------------------
% Package options (documented)
% ----------------------------------------------------------------
% Geometry scope:
%   apply-geometry            -> force A4 + zero margins while cover is built
%   no-apply-geometry         -> do not alter host geometry (default)
%
% Typesetting freeze:
%   global-typesetting        -> apply freeze globally (intrusive)
%   no-global-typesetting     -> no global freeze (default)
%
% Small-caps handling:
%   local-smallcaps-fix       -> \textsc -> \MakeUppercase only inside cover (default)
%   no-local-smallcaps-fix    -> disable local fix
%   global-smallcaps-fix      -> global override (intrusive)
%   no-global-smallcaps-fix   -> disable global fix (default)
%
% Main font handling:
%   global-mainfont           -> set Times New Roman globally (default)
%   local-mainfont            -> set Times only during cover rendering
%   no-mainfont-override      -> keep host document main font
%
% SVG support:
%   svg-support               -> enable SVG logos (default, needs shell-escape)
%   no-svg-support            -> disable SVG path handling (pdf/png only)
%
% Debug overlay:
%   debug                     -> show mm grid + bounding boxes on cover
%   no-debug                  -> hide debug overlay (default)
% ----------------------------------------------------------------
\usepackage[
  no-apply-geometry,
  no-global-typesetting,
  global-smallcaps-fix,
  local-smallcaps-fix,
  global-mainfont,
  svg-support,
  no-debug
]{ihedn-cover}

% ----------------------------------------------------------------
% Bibliography stack (BibLaTeX/Biber, style from subtree vendor/)
% ----------------------------------------------------------------
\usepackage{polyglossia}
\setdefaultlanguage{french}
\makeatletter
\g@addto@macro\captionsfrench{%
  \renewcommand{\tablename}{Tableau}%
  \renewcommand{\figurename}{Figure}%
}
\makeatother
\usepackage{csquotes}
\usepackage{xurl}
\usepackage{hyperref}
\makeatletter
\let\ihedn@orig@input@path\input@path
\def\input@path{{vendor/biblatex-sciences-po-fr/style/}}
\makeatother
\usepackage[
  backend=biber,
  style=sciencespo,
  hyperref=true,
  autolang=other
]{biblatex}
\AtBeginBibliography{\sloppy\setlength{\emergencystretch}{3em}}
\makeatletter
\let\input@path\ihedn@orig@input@path
\makeatother
\addbibresource{bibliographies/bib/rapport.bib}

% ----------------------------------------------------------------
% tcolorbox
% ----------------------------------------------------------------
\usepackage[most]{tcolorbox}
\tcbuselibrary{breakable}

% Example: use Arial for the full report + cover.
% Uncomment and switch package option to no-mainfont-override.
% \setmainfont{Arial}[
%   UprightFont    = *,
%   BoldFont       = * Bold,
%   ItalicFont     = * Italic,
%   BoldItalicFont = * Bold Italic
% ]

%==================================================
% Liste des propositions
%==================================================

\makeatletter

\newcommand{\listpropositionname}{Liste des propositions}

\newcommand{\listofpropositions}{
  \section*{\listpropositionname}
  \addcontentsline{toc}{section}{\listpropositionname}
  \@starttoc{lop}
}

\newcommand*\l@proposition{\@dottedtocline{1}{1.5em}{2.3em}}

\makeatother

%==================================================
% Compteur des propositions
%==================================================

\newcounter{proposition}

%==================================================
% Style graphique des propositions
%==================================================

\newtcolorbox{propositionbox}[2][]{
  lower separated=false,
  colback=white!80!gray,
  colframe=white!20!black,
  fonttitle=\bfseries,
  colbacktitle=white!30!gray,
  coltitle=black,
  enhanced,
  sharp corners,
  attach boxed title to top left={
    xshift=0.5cm,
    yshift=-2mm
  },
  title=#2,
  #1
}

%==================================================
% Environnement proposition
%==================================================

\newenvironment{proposition}[1]
{
  \refstepcounter{proposition}

  \addcontentsline{lop}{proposition}{%
    \protect\numberline{\theproposition}#1%
  }

  \begin{propositionbox}
  {Proposition~\theproposition~: #1}
}
{
  \end{propositionbox}
}

% Renommage Liste des figures en Table des figures
\addto\captionsfrench{%
  \renewcommand{\listfigurename}{Table des figures}
}

\newcommand{\heure}[2]{{\NoAutoSpacing #1:#2}}

\DeclareFieldFormat[misc]{title}{#1}

% Activer ce package avec l'option none retire la césure.
%\usepackage[none]{hyphenat}

\usepackage{adjustbox}
\usepackage{pdfpages}

\makeatletter
\newcommand{\IhednSetPdfPageCount}[2]{%
  \begingroup
  \edef\ihedn@pdfpath{#2}%
  \ifXeTeX
    \global#1=\XeTeXpdfpagecount "\ihedn@pdfpath" %
  \else
    \ifLuaTeX
      \saveimageresource{\ihedn@pdfpath}%
      \global#1=\lastsavedimageresourcepages\relax
    \else
      \pdfximage{\ihedn@pdfpath}%
      \global#1=\pdflastximagepages\relax
    \fi
  \fi
  \endgroup
}
\makeatother

\newcount\IhednIncludedPdfPageCount
\newcommand{\IhednRenderPdfPagesInBoxes}[1]{%
  \IhednSetPdfPageCount{\IhednIncludedPdfPageCount}{#1}%
  \foreach \IhednIncludedPdfPage in {1,...,\the\IhednIncludedPdfPageCount}{%
    \begin{tcolorbox}[
      colback=gray!5,
      colframe=black,
      boxrule=0.5pt,
      arc=3mm
    ]
    \centering
    \includegraphics[page=\IhednIncludedPdfPage,width=0.95\textwidth]{#1}
    \end{tcolorbox}
    \ifnum\IhednIncludedPdfPage<\IhednIncludedPdfPageCount
      \clearpage
    \fi
  }%
}

\usepackage{tablefootnote}


\begin{document}

% ----------------------------------------------------------------
% Cover keys (all available keys are shown below)
% ----------------------------------------------------------------
% Header block (top-right), logo, title, report lines, subtitle, committee.
\PageDeGardeIHEDN{
    header_1={Union-IHEDN},
    header_2={Association des auditeurs IHEDN},
    header_3={région Paris / Île de France},
    date={le 15 février 2026},
    logo_path={Figures/logos_IHEDN/Logo_AR_16_PARIS_IDF.jpg},
    title={Crises majeures : quel rôle pour les citoyens engagés et les associations IHEDN ?},
    title_max_lines=2,
    title_fontsize={26.000pt},
    title_baselineskip={28.667pt},
    reportline_1={Rapport du Comité d'étude de l'Association régionale des auditeurs},
    reportline_2={de l'IHEDN région Paris / Île de France (AR 16)},
    subtitle={Compte rendu de la réunion du 12 février 2026},
    committee={Comité d'étude, d'octobre 2025 à juin 2026, composé de : Mme Valérie Arlé-Faivre, \textit{présidente} ; M. Aurélien Duvignac-Rosa, \textit{rapporteur} ; M. Pascal Bayot Duponchel ; M. Jean-Jacques Cleynen ; Mme Valérie Cowan ; M. Alain Fontan ; M. Pierre-Marie Giard ; Mme Muriel Joyeux ; M. Bruno Leray ; M. Yves-Henri Renhas ; M. Marc Roure ; Mme Isabelle de Segonzac ; M. Gérard Turck ; Mme Marie Danielle Vasquez-Duchêne.}
}

\section*{Orientation générale des travaux}

La séance du 12 février 2026 s'inscrit dans la continuité des réflexions engagées sur la mobilisation des auditeurs IHEDN de l'AR 16  dans la gestion des crises majeures.

Le point central des échanges a porté sur la \textbf{nécessité de mieux connaître les membres de l'association}, par l'intermédiaire d'une \textbf{enquête rigoureuse, crédible et structurée}.

L'envoi d'un questionnaire constitue en soi un acte de communication : il manifeste une attention portée aux auditeurs et contribue à raviver le lien entre les membres.

Au cours de cette réunion, Jean-Jacques Cleynen a présenté ses premiers éléments de recherche qui constitueront la base rédactionnelle du rapport de comité.

Les éléments développés dans ce compte rendu en sont un résumé étoffé des divers échanges qui ont eu lieu le 12 Février 2026.

\section*{Temporalité de l'engagement des auditeurs de l'AR 16}

Une clarification importante a été apportée concernant la \textbf{temporalité de l'engagement des auditeurs} :

\begin{itemize}
    \item \textbf{Avant la crise} : rôle pertinent en matière de préparation, de sensibilisation et de diffusion d'une culture de résilience ;
    \item \textbf{Pendant la crise} : l'organisation étant strictement hiérarchisée, les auditeurs se mettent prioritairement au service des administrations ou corps auxquels ils appartiennent ;
    \item \textbf{Après la crise} : participation possible aux retours d'expérience (RETEX) et à l'analyse sociologique ou organisationnelle.
\end{itemize}

L'AR 16 peut fédérer ses membres, mais chacun conserve son rattachement institutionnel principal. Ainsi, la place qu'ils occupent au sein de la crise est un axe de réflexion majeur.

\section*{Positionnement dans la hiérarchie de crise}

La question centrale demeure : \textbf{comment intégrer l'AR 16 dans la hiérarchie de gestion de crise ?}

Les références opérationnelles suivantes peuvent constituer un axe de réflexion :
\begin{itemize}
    \item Plan ORSEC (Organisation de la Réponse de SEcurité Civile)\footnote{\href{https://www.hauts-de-seine.gouv.fr/Actions-de-l-Etat/Securite-et-Defense/Securite-civile/Le-Plan-ORSEC}{Le Plan ORSEC} [consulté le 16 févr. 2026].};
    \item PCS (Plan communal de sauvegarde) et PCIS (Plans intercommunaux de sauvegarde)\footnote{\href{https://www.securite-civile.interieur.gouv.fr/reagir/comment-se-preparer-face-aux-risques/plans-communaux-et}{Les plans communaux et intercommunaux de sauvegarde (PCS / PICS)} [consulté le 16 févr. 2026].} ;
    \item Rôle du DMD (Délégué Militaire Départemental)\footnote{\href{https://www.eure.gouv.fr/Services-de-l-Etat/Defense-memoire-anciens-combattants/Delegation-militaire-departementale-DMD2/MISSIONS}{Délégation militaire départementale (DMD) -- Missions} [consulté le 16 févr. 2026].}.
\end{itemize}

%\footnote{\href{https://www.securite-civile.interieur.gouv.fr/actualites/actualites/remise-du-rapport-de-synthese-du-beauvau-de-securite-civile-aux-ministres}{Présentation par le ministre du rapport de synthèse du Beauvau de la Sécurité civile }}

L'exploitation des résultats du questionnaire pourrait constituer un socle pour proposer un positionnement crédible de l'AR 16, notamment en amont et en aval des crises.
Cette démarche permettrait également de définir les profils disponibles et volontaires.

\section*{Fédérer et cartographier les compétences}

Le questionnaire aurait plusieurs objectifs :
\begin{itemize}
    \item Identifier le niveau d'implication souhaité en cas de crise ;
    \item Distinguer actifs et retraités ;
    \item Identifier, parmi les membres, ceux réellement volontaires ;
    \item Cartographier les compétences afin de définir la nature des crises auxquelles les auditeurs pourraient contribuer.
\end{itemize}

La question de la mise à jour de l'annuaire IHEDN a été soulevée. L'annuaire 2025 apparaît lacunaire.  
Une actualisation minimale (nom, fonction, compétences principales) constituerait un préalable structurant.

En complément de ces éléments, la nature de la crise majeure à laquelle serait confronté le pays est également un axe majeur de réflexion.

\section*{Définition et nature de la crise majeure}

Un besoin de clarification conceptuelle a été exprimé : 
\begin{itemize}
    \item Quelle est la définition exacte d'une crise majeure ?
    \item Quelles natures de crises justifient ou permettent l'implication des membres de l'AR 16 ?
\end{itemize}

Pour rappel, la gestion des crises majeures (catastrophes naturelles, attentats, accidents industriels, pandémie, etc.) en France repose sur un ensemble de procédures reglementées et une organisation hiérarchisée entre les territoires, préfectures et mairies.

Une recherche bibliographique pourrait être engagée afin de consolider ce cadre théorique.

A ce titre, l'expérience du COVID constituerait un terrain d'étude pertinent pour analyser la sociologie de la gestion de crise.

Le type de crise établi, des solutions opérationnelles doivent être envisagées.

\section*{Propositions opérationnelles}

L'implication de l'AR 16 dans le cadre de la gestion de crise majeur pourrait se materialiser de différentes manières:

\begin{itemize}
    \item Développer des préparations à la gestion de crise à destination des enfants (vulgarisation et acculturation) ;
    \item Transformer les bonnes pratiques en habitudes collectives ;
    \item Créer un lien entre les référents défense municipaux et les membres de l'AR 16 de leur commune ;
    \item Mettre en place un système d'alerte (ex. SMS) permettant de mutualiser et d'activer rapidement le réseau, en articulation avec le DMD ;
    \item Raviver le lien entre membres pour favoriser la diversité des profils et des expertises.
\end{itemize}

\newpage

\section*{Conclusion}

La réunion a confirmé la pertinence d'une démarche structurée visant à :

\begin{itemize}
    \item Mieux connaître les membres de l'AR 16 ;
    \item Clarifier leur positionnement avant et après les crises ;
    \item Proposer une intégration crédible dans les dispositifs existants ;
    \item Renforcer la cohésion et la lisibilité de l'association.
\end{itemize}

Le 20 Février 2026 à \heure{16}{00}, un échange avec un sociologue de l'IRSEM et des membres du comité a eu lieu pour enrichir la réflexion.

Les éléments issus de cet échange seront transmis ultérieurement.

\end{document}